% Generated draft: requires mathematical and source review.
\documentclass[11pt]{article}
\usepackage[margin=25mm]{geometry}
\usepackage{fontspec}
\setmainfont{Latin Modern Roman}
\usepackage{amsmath,amssymb,mathtools,unicode-math}
\setmathfont{Latin Modern Math}
\usepackage{enumitem,needspace,xurl,hyperref}
\setlength{\parindent}{0pt}
\setlength{\parskip}{5pt}
\setlength{\emergencystretch}{3em}
\newcommand{\N}{\mathbb{N}}
\newcommand{\Z}{\mathbb{Z}}
\newcommand{\Q}{\mathbb{Q}}
\newcommand{\R}{\mathbb{R}}
\newcommand{\C}{\mathbb{C}}
\newcommand{\F}{\mathbb{F}}
\newcommand{\A}{\mathbb{A}}
\newcommand{\PP}{\mathbb{P}}
\newcommand{\E}{\mathbb{E}}
\newcommand{\eps}{\varepsilon}
\newcommand{\dd}{\,\mathrm{d}}
\begin{document}
\section*{UnsolvedMath ID: 30000300}
\textbf{Source problem code:} catalogue entry
\subsection{Definitions and mathematical statement}
In the usual compact complex setting, let $(M,J)$ be a compact complex manifold and let $\Omega\in H^{1,1}(M,\R)$ be a K\"ahler class. Its representatives are positive real closed $(1,1)$-forms $\omega$ with $[\omega]=\Omega$. Such a form determines a K\"ahler metric and scalar curvature $S_\omega$.

A metric is extremal when its complex scalar-curvature gradient is a holomorphic vector field:
\[
 \bar\partial\bigl(\operatorname{grad}^{1,0}_{\omega}S_\omega\bigr)=0.
\]
Equivalently in the compact setting it is a critical point of the Calabi scalar-curvature functional within the K\"ahler class. The target is to decide whether the given class contains such a representative, and, if one exists, in what sense it is unique.

The standard geometrically meaningful uniqueness question is modulo holomorphic automorphisms preserving the class. Put
\[
 \operatorname{Aut}(M,\Omega)=\{F\in\operatorname{Aut}(M):F^*\Omega=\Omega\}.
\]
For extremal $\omega_1,\omega_2\in\Omega$, this asks for $F$ in the relevant automorphism group with $F^*\omega_2=\omega_1$. The source simply says ``unique'' and does not specify this quotient or the identity-component convention. Literal equality of representatives would be stronger and is not invariant under nontrivial class-preserving automorphisms, so that interpretation is not silently imposed.

The extract also does not explicitly state compactness or additional conditions selecting classes expected to admit extremal metrics. The displayed compact formulation is conventional; the source question is not an assertion that every arbitrary K\"ahler class automatically has a solution. No integrality, projectivity or constant-scalar-curvature condition is added, since extremal metrics may have nonconstant scalar curvature.

The literature assessment checked on 21 August 2026 reports existence and uniqueness in important stable and geometric settings while the general extremal Yau--Tian--Donaldson existence problem remains open in that assessment. These are dated partial-progress claims, not an independently updated classification or a statement of literal unquotiented uniqueness.
\subsection{Short English statement}
For a given K\"ahler class, determine whether an extremal metric exists and whether it is unique in the appropriate sense, normally modulo holomorphic automorphisms. The source does not supply a universal existence criterion or a precise uniqueness quotient.
\subsection{Sources}
\begin{itemize}
\item \emph{Analysis and Geometric Singularities} (2006), Oberwolfach Reports, DOI 10.4171/owr/2005/39, pp. 2187--2240, \url{https://doi.org/10.4171/owr/2005/39}; proposed year supplied as 2005.
\item Literature assessment checked 21 August 2026: original report, \url{https://doi.org/10.4171/owr/2005/39}; ``Extremal Kahler metrics,'' \url{https://arxiv.org/abs/1405.4836}.
\item Local UnsolvedMath registry, permanent ID $30000300$, exact source code \texttt{catalogue entry}, created 25 August 2026 and updated 16 September 2026. Corrected-and-verified status records K\"ahler notation and punctuation normalization while preserving both existence and uniqueness questions.
\end{itemize}
\end{document}
